% WHAT A WORD HYPHENATES TO WHEN THE HYPHEN CHARACTER IS NOT THERE.
%
% Eight lines that draw a real divergence, found by cutting trip's first pass
% down until it was the only thing left. In trip's own font, with \hyphenchar
% set to a character the font has not got, hyphenating `111':
%
%   ref:    Missing character: There is no ^^A in font trip!
%           Missing character: There is no ^^A in font trip!
%           ..\discretionary replacing 70
%           ..\discretionary replacing 2
%
%   hstex:  ..\discretionary replacing 5
%           ..\discretionary replacing 67
%
% TWO FAULTS, AND THEY ARE SEPARATE. [B] separates them: it asks for the
% same word with `\hyphenchar' set to `-', which the font DOES have.
%
%   - The missing hyphen goes unnamed. In [A] the reference draws
%     `Missing character' twice and this engine not at all. In [B], where
%     there is no missing character, neither draws one.
%   - The discretionary's REPLACE COUNT differs, in [A] and [B] alike: 70
%     against 5, and 2 against 67. The ligature nodes on either side of it
%     agree exactly -- both engines draw `\trip 5 (ligature |1)' and
%     `\trip 7 (ligature |)' -- so the two agree about what the word is made
%     of and disagree about how much of it the discretionary puts back.
%
% So the second fault is not about the hyphen at all. It is about what a
% discretionary made while a LIGATURE is being reconstituted replaces.
%
% WHY IT IS THE TRIP FONT AND NOT THE HYPHENCHAR. With cmr10 the two agree
% exactly at every \hyphenchar tried: 200, which cmr10 has not got, draws
% four warnings from both; 1 and 5, which it has, draw none from either.
% trip.tfm says character 1 has width index nought, which is a TFM's way of
% saying the character is not there, and this engine's own reader honours
% that -- it leaves the metric's tag at -1 and make_hyphen_node warns on a
% tag below nought. So the warning is not missing because the character is
% thought to be present; it is missing because make_hyphen_node is never
% reached for this word.
%
% WHAT IS DIFFERENT ABOUT THIS WORD is that trip.tfm gives character 49 --
% the `1' -- a ligature program (tag 1, remainder 9), and the word is three
% of them. `replacing 70' says the reference puts seventy nodes back where
% this engine puts back five, which is the size of what trip's ligature
% program builds out of them.
%
% WHERE THIS CAME FROM. trip's first pass has one divergence left, a single
% badness -- see what-a-break-is-weighed-as-in-trips-first-pass.tex -- and
% bisecting it says it needs `\righthyphenmin0', a language change part way
% through the paragraph, and \lccode`1 set so that `1' is a letter at all.
% This file is what is left when those are cut away too. It is NOT known to
% be the same fault: trip's own pass draws the same material either way and
% differs only in the trace, while this draws different discretionaries. It
% is recorded because it is small, certain, and in the same machinery.
\catcode`\{=1 \catcode`\}=2
\tracingonline=1 \tracinglostchars=2 \showboxbreadth=99 \showboxdepth=99
\font\trip=trip \trip
\hsize=1pt \parindent=0pt \pretolerance=-1 \tolerance=10000 \hbadness=10000
\lccode`1=`1 \patterns{111}
\message{[A trip, with a hyphenchar the font has not got]}
\hyphenchar\trip=1
\setbox0=\vbox{\noindent 1 111\par}\showbox0
\message{[B the same word with a hyphenchar the font does have]}
\hyphenchar\trip=`-
\setbox0=\vbox{\noindent 1 111\par}\showbox0
\message{[C every character below a space, which used to be drawn raw]}
% FIXED. show_ascii let 9, 10 and 11 through unescaped, because the same
% predicate governs what the engine writes of its own text -- where a newline
% must stay a newline. A character out of a document is not that: the
% reference draws 8 to 12 alike as ^^H ^^I ^^J ^^K ^^L, and cmr10's `ff'
% ligature IS character 11, so every box holding one came out with a raw
% vertical tab in it. Found here, and nothing to do with hyphenation.
\newlinechar=-1
\font\tenrm=cmr10 \tenrm
\setbox0=\hbox{\char8 \char9 \char10 \char11 \char12 }\showbox0
\message{[D an ff ligature, which is character 11]}
\lccode`o=`o \lccode`f=`f \lccode`e=`e \lccode`r=`r
\patterns{o1ffer}
\setbox0=\vbox{\noindent x offer\par}\showbox0
\message{[E and cmr10 hyphenchars, which agree everywhere]}
\lccode`a=`a \patterns{a1a}
\hyphenchar\tenrm=200 \setbox0=\vbox{\noindent x aaaaa\par}
\hyphenchar\tenrm=1 \setbox0=\vbox{\noindent x aaaaa\par}
\hyphenchar\tenrm=5 \setbox0=\vbox{\noindent x aaaaa\par}
\message{[done]}
\end
